% ~600 words. Three-beat argument: evidence gap = access gap; benchmark
% saturation; FLAIR's three contributions.
\section{Introduction}

The bottleneck in clinical decision support (CDS) for the intensive
care unit is no longer the supply of models. Predictive algorithms for
sepsis, acute kidney injury, extubation readiness, and ICU mortality
have been published at a rate the literature struggles to absorb. What
the field is missing is the evidence that any given model behaves
acceptably outside the institution that built it. The body of evidence
required to support routine clinical use has not kept pace with the
proliferation of models, and that gap is, in our view, primarily an
\emph{access} gap: representative, multi-institutional data on which
to test a model rarely sit within reach of the team that needs them.

\paragraph{The evidence gap is an access gap.}
External validation of even flagship EHR-based CDS has remained
stubbornly narrow. The Epic Sepsis Model, deployed in hundreds of
United States hospitals, was first independently evaluated at exactly
one site, with discrimination falling to AUC~0.63 against vendor-claimed
ranges of 0.76--0.83 \citep{wong2021external,habib2021jamaeditorial}.
Duke's Sepsis Watch, an exemplar academic deployment, took more than
six years to be assessed at its first external site
\citep{sendak2025sepsiswatchportability}. The Targeted Real-time Early
Warning System (TREWS) reached five sites, all within a single health
system \citep{adams2022trews}. Mayo Clinic's electrocardiogram-based
low-ejection-fraction screen, taken through the FDA pathway, required
multi-year coordination to cover roughly four to six external sites
\citep{carter2025mayoecgai}. Each of these efforts required a bilateral
data use agreement, shipping de-identified patient data to a
coordinating analytic center. None scale to dozens of sites at the pace
the rest of the AI ecosystem moves.

\paragraph{Public benchmarks saturate; private benchmarks don't scale.}
The community has filled some of the gap with open ICU benchmarks
\citep{vandewater2023yaib,yeche2021hirid} built on MIMIC, eICU, HiRID,
and AUMCdb. These tools are essential for development but inherit the
selection biases of the source datasets \citep{roosli2022mimicfairness},
and the literature now documents how repeated tuning against a fixed
public benchmark produces ``illusory generalizability''
\citep{recht2019imagenet,chekroud2024illusory}. Federation-by-data-use-
agreement helps but is narrow in scope: the Observational Health Data
Sciences and Informatics network federates SQL queries and coefficient
sharing for logistic models on a pre-harmonized common data model
\citep{reps2020ohdsistroke}; MedPerf and the FeTS Challenge prove that
a portal-orchestrated, model-to-data benchmarking ecosystem works at
scale, but for imaging \citep{karargyris2023medperf,pati2022fets}. None
of these combine modern, longitudinal EHR-based CDS, a clinically
meaningful ICU schema, and an extensible task catalog open to the
community.

\paragraph{Contribution.}
We present FLAIR (Federated Learning Assessment for ICU Research), a
benchmark framework with three load-bearing properties. First, it
operates on the \emph{model-to-data} paradigm pioneered by the DREAM
Challenges and applied at scale to EHRs in the EHR DREAM Challenge
\citep{guinney2018dream,bergquist2024ehrdream}, but extended from
one-time competitions to a standing, repeatable infrastructure.
Second, it is anchored on the Common Longitudinal ICU Format (CLIF), a
critical-care common data element schema adopted by 12+ academic
medical centers \citep{rojas2025clif}, conceptually parallel to mCODE
for oncology \citep{mcode2024hl7} and complementary to the Medical
Event Data Standard \citep{arnrich2025meds}. Third, FLAIR is released open source (Apache~2.0,
\texttt{github.com/clif-consortium/FLAIR}): the task catalog,
evaluation protocol, and submission contract are extensible by
GitHub pull request, with each new task added as a single Python
module under \texttt{flair\_benchmark/tasks/}. The benchmark
co-evolves with clinical need rather than ossifying around a single
author group's priors, and the five-task v1 suite below is
explicitly a seed for that catalog rather than a fixed scope.

We organize the rest of the manuscript around three commitments: the
CLIF-anchored \emph{datasets} that participating sites contribute
(\S\ref{sec:datasets}), the initial five-task \emph{benchmark} suite
together with the open contribution protocol for adding more tasks
(\S\ref{sec:benchmarks}), and the federated submission and evaluation
\emph{protocol} (\S\ref{sec:protocols}).
